% Appendix A, Chapter 2 — robustness checks. Content mirrors robustness_checks/README.md's index
% and the individual RCNN_note.md files. Regenerate/extend as further checks are run.

\section{Summary of checks run}

\begin{threeparttable}
\caption{Robustness checks, Chapter Two}
\begin{tabular}{p{1.1cm} p{3.5cm} p{4.3cm} p{3.6cm}}
\toprule
Check & Baseline & Alternative specification & Verdict \\
\midrule
RC01 & Four-bloc Pedersen index (Con/Lab/LD/Other) & Eight-party, fuller coding & \textbf{Survives}
--- break test identical, $r=0.99$ \\
RC02 & Figure 2.9 Panel B, level only & Adds persistence: rolling AR(1) of the composite &
\textbf{Survives, qualified} --- moderate persistence (0.36 whole-series), so not pure churn;
COVID-era exception flagged \\
RC03 & Figure 2.9 composite: $R^2 \geq 0.40$ inclusion floor, 10\% Chow trim, 2020--21 COVID
exclusion & Floor lowered to $R^2 \geq 0.30$, recovering Immigration, Housing and Pensions
(5 $\rightarrow$ 8 items) & \textbf{Survives; floor retained at 0.40} --- same break year (1985),
$F$ falls 83.17 $\rightarrow$ 51.44, still $p<0.0001$ \\
RC04 & Figure 2.9's OLS-residual ``unexpected'' measure (5 domains, break 1985, $F=83.17$) &
Restricted slope\,$=1$ z-score gap, all 13 domains, no exclusion & \textbf{Rejected as a main
measure}, reasons disclosed --- see below \\
RC10 & Figure 2.15 U-turns: raw annual count, 221 events & (i) per 100,000 characters of chapter;
(ii) high-confidence events only; (iii) narrowly specified commitments only &
\textbf{Moves, reverses, survives} --- (i) reverses the early--late comparison, both reported;
(ii) reverses, but is \textbf{inadmissible} (era-dependent filter); (iii) survives \\
RC11 & Figure 2.16 turbulence index: segmented (bend) fit, knot 2013 & Flat, linear trend,
quadratic, single break, thin-plate spline --- compared on BIC, AIC, leave-one-out CV and
AR(1)-error BIC & \textbf{Survives as a \emph{shape}; the form and date are not identified} ---
flat and linear rejected ($\Delta$BIC 27.4, 16.8); bend, break and spline indistinguishable
($\Delta$BIC $\leq0.71$) \\
RC12 & The volatility curse (Campello \& Zucco) as the bridge from turbulence to democratic
consequence & Attribution link measured two ways ($R^2$ and standardised slope) over three eras,
11 agenda items; plus four related causal claims & \textbf{Not found} --- $R^2$ decays but the
range-restriction-proof slope does not ($p=0.65$); four further claims all vanish under a time
trend \\
\bottomrule
\end{tabular}
\begin{tablenotes}\footnotesize
\item Full detail for each check is in its own folder: \texttt{robustness\_checks/RCNN\_chX\_
shortname/RCNN\_note.md}, with regenerated outputs alongside.
\item Two further checks exist as their own folders and are not yet written up here: a replication
of the Camargo analysis, and a single-mention variant of the public-opinion series.
\end{tablenotes}
\end{threeparttable}

\section{RC03 --- the inclusion floor, and why it was retained}

The Figure 2.9 composite includes issue domains whose tracking regression clears an $R^2$ floor of
0.40. That threshold is discretionary, so it was tested directly. Lowering it to 0.30 recovers three
further domains (Immigration, Housing, Pensions), taking the composite from five items to eight.

The finding does not depend on the choice: the break year is unchanged at 1985, and the statistic
weakens but remains overwhelming ($F$ from 83.17 to 51.44, $p<0.0001$ throughout).

The floor was nonetheless retained at 0.40, for a stated reason rather than by default. The
threshold sits inside a genuine gap in the data --- five domains cluster between $R^2$ of 0.60 and
0.90, and the next-best item sits at 0.35, so 0.40 separates two natural groups rather than cutting
through a continuum. The RC04 diagnosis below independently reinforced the decision, by showing what
happens when a looser rule admits an item that then dominates the result. Nothing in the live
pipeline was changed by this check.

Sensitivity to the other two discretionary parameters in this composite --- the 10\% Chow trim and
the 2020--21 COVID exclusion window --- has not yet been run, and remains outstanding.

\section{RC04 --- an alternative measure, and the diagnosis that rejected it}

This check is recorded in full because its interim result was attractive and wrong, and the
reasoning that rejected it is the substantive point.

The baseline ``unexpected movement'' measure is an OLS residual across five domains, breaking in
1985. The alternative --- a restricted slope\,$=1$ z-score gap, using all thirteen domains with no
sign exclusion --- produced a composite with a quite different and initially appealing shape: a
break in 2019 on a rising series, $F=23.61$.

Item-level diagnosis found this was almost entirely one series. Poverty shifted by 2.31, roughly
five times the next largest item, while six other domains moved in the opposite direction. A median
composite, which is robust to a single dominating item, removes the rise altogether: the break moves
to 2009, $F=8.93$, and \emph{both} segments decline.

The check therefore did more than reject a candidate measure. It independently vindicated the
sign-exclusion rule used in the main specification, which already screens out the one item that
distorted the alternative. Closed 30 July 2026; full working in \texttt{RC04\_note.md}.

\section{RC05 --- sensitivity of the policy-agenda and spending findings}

The two series in the policy-instability section (Figure~2.11 and the spending-composition figure)
each yield a concentration result and a change result. The concentration results are robust; the
\emph{dates} attached to them are not, and neither is the change result. This section sets out what
each finding does and does not survive, since that distinction is the whole of what can be claimed.

\begin{threeparttable}
\caption{Sensitivity of the policy-agenda and spending findings}
\begin{tabular}{p{2.6cm} p{3.6cm} p{4.2cm} p{2.6cm}}
\toprule
Finding & Baseline & Alternative specification & Verdict \\
\midrule
Agenda concentration, break year & Full series from 1972: peak 1997, $F=25.84$, $p<0.0001$ &
Restricted to 1979--2026, the book's window and the range the figure displays: $F=8.81$, $p=0.010$;
from 1980, $F=6.39$, $p=0.051$ & \textbf{Window-sensitive} --- the 1979 window is used throughout,
and the reported $F$ falls threefold \\
Agenda concentration, coder & CAP human coding, 1979--2012: peak 1997, $F=8.15$, $p=0.018$ &
This project's independent coding, same years: peak 1997, $F=6.35$, $p=0.059$ &
\textbf{Survives} --- both coders locate the same year on matched data \\
Agenda concentration, break date & Sweep peak 1992 & 90\% bootstrap interval 1987--2014, against an
admissible window of 1984--2017; $F$ exceeds the critical value at every candidate year &
\textbf{Not identified} --- no date is reported \\
Agenda churn & CAP coding, 1979--2012: $+0.00276$/yr, $p=0.0056$ & This project's coding, same
years: $+0.00135$/yr, $p=0.283$ & \textbf{Resolution effect}, see note \\
Spending concentration & $-0.0230$/yr, $p<0.0001$ & Excluding 2020--22: $-0.0263$/yr, $p<0.0001$;
CAP segment alone $-0.0562$, $p<0.0001$ & \textbf{Survives} --- but flat within the PESA segment
alone ($p=0.26$), so the decline is concentrated pre-2005 \\
Spending change & All years: $+0.0117$/yr, $p=0.675$ & Excluding 2020: $-0.0132$/yr, $p=0.548$;
excluding 2020--22: $-0.0406$/yr, $p=0.024$ & \textbf{Sign depends on one observation} --- no claim
of rising volatility is made \\
Spending break year & Sweep peak 2018, $F=11.62$, $p=0.026$ & Excluding 2020 the peak moves to 1991
and loses significance ($p=0.130$) & \textbf{Artefact of the pandemic year} \\
Agenda concentration, unit density & 2019--26 reads denser by ratio than 2000--17 & Rarefied to a
common 43 units: the gap widens (4.20 vs.\ 3.27 raw), $t=-7.67$, $p<0.0001$ &
\textbf{Survives} --- not a segmentation artefact \\
\bottomrule
\end{tabular}
\begin{tablenotes}\footnotesize
\item \textbf{The agenda churn discrepancy is a resolution effect, not a disagreement.} Subsampling
the human-coded series down to this project's unit counts reproduces this project's null almost
exactly ($+0.00095$/yr; 8\% of 200 bootstrap draws significant, against 38\% at a constant 46 units
and $p=0.0056$ unrarefied). The two codings agree on topics ($r=0.892$) and differ in how finely
they cut, so the human-coded series is the better-resolved one and is the one the trend should be
quoted from. This project's flat slope is a measurement limitation, not a competing estimate.
\item \textbf{Both break dates are unidentified, and both underlying changes are real.} Where a
structural-break test rejects a single-trend model at every candidate year, it is detecting a change
of shape rather than locating an event. The concentration trends are significant at $p<0.001$ in
both series and run in opposite directions --- the agenda broadens while spending concentrates.
\item \textbf{The spending series is dominated by two observations}, 1990--91 and 2020--21, and in
both cases by one function (Economic affairs, whose share falls 52\% then rises 95\%, and later rises
128\% then falls 45\%). Read the series as a response to two recessions rather than as a general
loosening of budgetary discipline.
\item Full detail in \texttt{policy\_instability/agenda\_recode/METHOD.md} \S4t--\S5e and in the
\texttt{diagnostics/} scripts, each of which regenerates the numbers above.
\end{tablenotes}
\end{threeparttable}

\section{RC06 --- the single-break specification, and why the start date matters}

The break tests throughout this chapter fit \emph{one} break. That specification is tested here
against the two series in the policy-instability section, and it does not survive. Both series are
curved rather than trending, both prefer more than one break on a formal criterion, and both give
materially different answers depending on where the series is allowed to begin. None of this is a
coding artefact --- it is what the data do.

\subsection*{The series are U-shaped, not trending}

Fitting a quadratic in place of a straight line raises $R^2$ from 0.255 to 0.487 for the government
agenda ($F=18.52$, $p=0.0001$) and from 0.553 to 0.817 for spending composition ($F=62.32$,
$p<0.00001$). Era means show the shape directly, and show that the 1980s are the high end of it
rather than a quiet baseline:

\begin{center}
\begin{tabular}{lrrrr}
\toprule
& 1980s & 1990s & 2000s & 2010s-- \\
\midrule
Agenda (effective topics) & 11.95 & 10.23 & 11.98 & 13.80 \\
Spending (effective functions) & 5.84 & 5.02 & 5.09 & 4.90 \\
\bottomrule
\end{tabular}
\end{center}

The 1980s--1990s difference is significant in both series ($p=0.049$ and $p<0.0001$). The
government agenda was as broad in the 1980s as in the 2000s; public spending was at its most
dispersed. The mid-1990s are the trough of both, not the floor from which a rise begins.

\subsection*{The data prefer more than one break}

Selecting the number of breaks by BIC (\texttt{strucchange::breakpoints}, trend model, 15\% trim)
returns \textbf{two} breaks for the agenda (1992 and 2003) and \textbf{four} for spending (1987,
1996, 2007 and 2013). A single-break test applied to a series with two or four turning points is
misspecified, and this explains a pattern reported elsewhere in these appendices: the $F$-curve
never falls below its critical value, and the bootstrap interval for the break year spans almost the
whole admissible window. A wandering series rejects the single-trend model at every candidate year,
and the ``break'' it reports is wherever the best single split happens to land.

\subsection*{The start date changes the finding}

\begin{threeparttable}
\caption{Trend in concentration by start year}
\begin{tabular}{lrrrr}
\toprule
Series starts & \multicolumn{2}{c}{Government agenda} & \multicolumn{2}{c}{Public spending} \\
& slope/yr & $p$ & slope/yr & $p$ \\
\midrule
1979 & $+0.0803$ & 0.0005 & $-0.0230$ & $<0.0001$ \\
1990 & $+0.1688$ & $<0.0001$ & $-0.0045$ & 0.163 \\
1997 & $+0.1545$ & 0.0001 & $+0.0012$ & 0.777 \\
1999 & $+0.1529$ & 0.0006 & $-0.0014$ & 0.773 \\
\bottomrule
\end{tabular}
\begin{tablenotes}\footnotesize
\item \textbf{This is the strongest single caution in these appendices.} The agenda series reaches
its minimum in 1993; a start date shortly after it roughly doubles the apparent rate of broadening.
The spending concentration result --- the largest $F$ statistic anywhere in this section --- is a
phenomenon of 1979 to the late 1990s and \textbf{disappears entirely} from a 1997 start.
\item A start year chosen for its historical tidiness can therefore manufacture one finding and
erase another. 1997 is attractive as a baseline on institutional grounds and is, for these two
series, close to the worst available choice on statistical ones: it sits just after the trough of
one series and after the whole of the movement in the other.
\item The window used throughout is 1979--2026 for the agenda and 1979--2024 for spending, chosen
before these tests were run and retained because it is the book's frame, not because it is
favourable. Readers should treat the \emph{shape} of both series as the finding and any single
slope or break year as conditional on the window.
\item Reproduced by \texttt{policy\_instability/diagnostics/eighties.R}.
\end{tablenotes}
\end{threeparttable}

\section{RC07 --- the National Audit Office delivery-failure series}

Figure~2.13 rests on a machine-coded reading of every National Audit Office value-for-money
report, 1984--2025. The August 2026 re-code changed both the extract the coder reads and the
severity scale it applies, and moved the finding materially. This section reports what moved,
and plots every specification considered and rejected.

\subsection*{Why the series was re-coded}

The original extract took the first passage of continuous prose in each report. In National Audit
Office reports that passage is scene-setting --- what the body does, how large the programme is ---
and the verdict sits deeper. Hand-coding of eleven reports, blind to the machine's answer,
established the cost: an expert reader could code only two of five reports from the original
extract. Re-extracting from the report's own findings section lifted that to eight of eleven.
The re-code raised the measured failure rate in every period after 1996 by seven to eight percentage
points, and left the 1984--96 period essentially unchanged --- the defect was in modern reports,
not old ones, because older reports carry a clean summary heading and newer ones do not.

\subsection*{What is plotted, and what is not}

The figure plots a \emph{count}: serious delivery failures identified per year, estimated from
extracts carrying an explicit verdict and applied to all value-for-money reports published. The
proportion measure --- the more intuitive quantity, and the one shown in earlier presentations of
this work --- is not used, for the reason set out below.

\begin{figure}[htbp]\centering
\includegraphics[width=0.98\textwidth]{robustness_chapters/figures_ch02/fig_a2_14_nao_specifications_not_chosen.png}
\caption{National Audit Office delivery failure: the specifications not chosen. Panel A is the
superseded series; B and C are proportion measures on the corrected extract; D restricts to failure
attributed to the audited body's own decisions; E is the count without the verdict-bearing
restriction; F is the extract coverage that disqualifies a raw count within that subset.}
\label{fig:ch02naoalt}
\end{figure}

\subsection*{Publication volume: the confound that decided the measure}

National Audit Office output rose to roughly 75 reports a year in 2005--13 and fell to about 57 in
2014--25. The rise in the \emph{share} of reports finding serious failure coincides with that
decline, and a referee's objection follows directly: publish fewer reports, more selectively, and
the share rises without any change in government performance. The objection has force. Conditional
on volume, the post-2015 excess in the proportion measure is not statistically significant
(Poisson with a volume control and a year trend: $+0.203$, $p=0.258$), and adding volume as a
regressor removes the break under Bai--Perron.

The count measure is reported instead because it inverts the objection rather than evading it:
from 2005--13 to 2014--25 the National Audit Office published about a quarter \emph{fewer} reports
and identified about a third \emph{more} serious failures. That contrast requires no model. It does
not dispose of the concern --- the Office may have changed what it chooses to audit, which these
data cannot test --- and readers should treat the break year as conditional on that.

\subsection*{Attribution: what cannot be claimed}

The re-code coded, separately from severity, whether each failure lay within the audited body's own
control. Externally-driven failure is far commoner in the recent period (21.3 per cent of
2014--25 reports against 8.1, 10.2 and 9.3 per cent in the three earlier periods), and restricting to
government-attributable failure cuts the rise by about a quarter.

On that restricted series \textbf{no statistically significant break survives} ($p=0.34$;
Bai--Perron selects zero breaks). The period contrast remains substantial --- 17.1 to 28.2 per cent
--- but the annual series is too noisy for the test once the numerator is restricted. This is a
limit on what may be claimed rather than evidence that nothing changed: \emph{a difference in period
means} is supportable, \emph{a structural break in government-attributable failure} is not.

\subsection*{Reproduction}

Series and both figures: \texttt{nao\_extension/diagnostics/build\_fig214\_v3.R} and
\texttt{build\_fig214\_robustness\_appendix.R}. Full method, including the coding prompts and the
instrument defects found and fixed, in \texttt{nao\_extension/METHOD.md}.

\section{RC08 --- the 1997--98 gap: recovery, and a substitute source rejected}

Figure 2.13 was built with 1997 and 1998 absent. Both years are now present, on the same
instrument as every other year. This section records how the gap arose, how it was closed, and
--- because the route that failed is more informative than the one that worked --- what was
tested and rejected along the way.

\subsection*{Why the years were missing}

The reports were never unpublished. The National Audit Office catalogues 41 reports for 1997 and
39 for 1998 --- 40 and 38 of them value-for-money studies --- each with a live page on its current
website. What failed was
retrieval. Every such page hands the reader to the UK Government Web Archive, and two conditions
must both hold before that archive returns a document rather than its own viewer page: the URL
requires the \texttt{id\_} modifier immediately after the fourteen-digit timestamp, and the
timestamp must be one under which the file was actually captured. The 2022 capture that the
Office's own pages link to is empty for these files; the 2017 capture holds them. Requests that
satisfy only the first condition return a 17.5KB HTML page which is easily mistaken for a short
document.

Applying both conditions recovered \textbf{69 of 69} outstanding reports as original PDFs --- 30 for
1997 and 39 for 1998. With the reports already held, 1997 now has 40 coded value-for-money reports
(21 carrying an explicit verdict) and 1998 has 39 (23). Both clear the inclusion thresholds of
RC03 unchanged, on original documents, using the same prompts as the rest of the corpus. No
substitute source enters the series and no caveat attaches to these two years on grounds of
provenance.

The corpus is not, however, on a single extractor, and an earlier version of this section wrongly
implied that it was. Three text routes coexist and are recorded per row: the v3 findings
extractor; a v2 first-prose fallback where no findings anchor could be located; and optical
character recognition for reports digitised as images. Their extent is given below.

\subsection*{A substitute source, tested and rejected}

Before the recovery succeeded, the archived National Audit Office press notices were assessed as a
replacement. They survive in quantity for both years (49 and 46), are written by the Office, and
state the Comptroller and Auditor General's verdict in plain terms. They were rejected on evidence
rather than principle, and the tests are reported here because the same reasoning governs any
future use of a substitute instrument.

\emph{Measurement.} For 297 reports the corpus holds both a press notice and the report itself.
Coded with the identical prompt, the press notices returned a serious-failure rate of 0.141
against 0.212 from the reports --- a deficit of 7.1 percentage points. Against an
independent-samples two-standard-error band this is 6.2 points; because the two codings are
\emph{paired} on the same 297 reports the appropriate band is 4.6, so the deficit is more clearly
established than the independent-samples figure suggests. The disagreement is not symmetric noise: of 63 reports the full text codes as serious
failure, the press notice agrees on 28 and records 30 as the milder category. The notices report
the criticism but omit the quantified detail --- the percentage overrun, the share of a portfolio,
the exact audit-opinion wording --- on which the severity threshold turns. A one-directional
error of this kind does not cancel in a count.

\emph{Correction attempted.} Because the fault is threshold placement rather than comprehension,
a few-shot correction was tried, with exemplars drawn only from pairs excluded from the test set
and a pass condition fixed in advance. It failed: the deficit closed from $-0.068$ to $-0.064$
against a two-standard-error band of 0.063, and $\kappa$ fell from 0.363 to 0.337.

\emph{Selection --- claimed, then withdrawn.} An earlier version of this section reported that
reports receiving a press notice were more likely to be serious than those that did not, and drew
on that to argue that measurement and selection biases ran in opposite directions. \textbf{That
statistic has been withdrawn.} Its comparison group was constructed by title matching, and for
1999--2008 the corpus holds 588 value-for-money reports against 561 archived notices: at most a
few dozen reports can genuinely lack a notice, yet the ``without'' group contained several
hundred. It was therefore composed overwhelmingly of matching failures rather than of reports
without a notice --- and since matching succeeds more readily on distinctive, prominent titles,
and prominent reports are likelier to be serious, the estimate was biased in the direction claimed
by construction. A selection effect may well exist; this analysis did not measure it, and would
need matching on report number rather than title. The measurement finding above stands on its own
and is sufficient to reject the substitute.

\emph{Human check.} Ten press notices drawn at random from 1997--98 were coded by hand, blind to
the model's codes, which were sealed in advance. Agreement was 6 of 10 --- reported here as a
qualitative check and not as a reliability coefficient, since at that sample size the interval is
too wide to support one. The two disagreements at the severe end were interpretable rather than
random: one was the qualified-audit-opinion trigger applied literally to a qualification affecting
one of five accounts, the other an attribution question where Community-wide losses were assigned
to a United Kingdom body. Both are recorded in the coding notes.

\subsection*{An inclusion threshold briefly relaxed, and restored}

While the years were still short of documents, the inclusion thresholds were lowered from 15
reports and 10 verdict-bearing extracts to 13 and 6. The change was applied to every year rather
than to 1997 alone --- it also admitted 1987, 1988 and 1999 --- and it left the break at 2015,
moving the sup-$F$ from 17.06 to 15.23. It has been \textbf{reverted}. It was a response to
missing documents, and the documents are no longer missing; the thresholds of RC03 now apply
unchanged throughout. It is recorded here because a threshold moved to admit a wanted year is a
specification chosen to obtain a result, and the reader is entitled to know that the possibility
arose and was closed.

\subsection*{What the completed series shows}

On 41 years the break remains \textbf{2015}, with sup-$F$ 16.13 and a permutation
$p$ of 0.0002. Leave-one-out returns 2015 in 40 of 41 runs, and Bai--Perron
continues to select a single break at 2015. Period means are 7.3, 14.9, 17.6 and 24.1
serious failures a year across 1984--96, 1997--2004, 2005--13 and 2014--25 --- a rise of
3.29 times from the first period to the last.

The volume qualification of RC07 survives the completed series and should continue to be stated:
conditional on publication volume and a linear time trend, the post-2015 excess is not
statistically significant ($+0.165$, $p=0.315$). That qualification is, however,
\textbf{specification-sensitive}. Removing the time trend while retaining the volume control makes
the post-2015 term significant ($+0.632$, $p<0.001$). This
is not a robustness gain: without a trend the indicator absorbs the general upward drift of the
series. Both specifications are printed by the supporting script so that the sensitivity is
visible rather than resting on which one happened to be run.

\subsection*{Coverage across the rest of the series}

Because the gap was caused by a retrieval fault rather than by absent documents, every year was
audited against the Office's own catalogue. Coverage runs at 96--100 per cent for 1984--99. The
apparent shortfall in later years is accounted for by document type rather than by missing
reports: the catalogue also contains financial audit certificates, impact case studies and the
Public Service Agreement data-system reviews, none of which is a value-for-money report and none
of which enters the corpus.

\subsection*{Reproduction}

Recovery and coding: \texttt{nao\_extension/pdf\_era/code\_ukgwa\_recovered.py}. Substitute-source
tests: \texttt{calibrate\_press\_notices.py}, \texttt{tune\_press\_notice\_coder.py} and
\texttt{selection\_test\_9798.py}, with the hand codes in \texttt{handcode\_pj.csv} against the
sealed model codes in \texttt{handcode\_model\_codes\_SEALED.csv}. Supporting tests:
\texttt{diagnostics/fig214\_supporting\_tests.R}. Series thresholds and their history:
\texttt{pdf\_era/rebuild\_series.py}. Full narrative, including three superseded accounts of the
gap and their disproof, in \texttt{nao\_extension/METHOD.md}.

\subsection*{Where the break falls, by specification}

The break date is the interpretive claim this section rests on, so it is worth setting out how it
behaves under every alternative tried. Table~\ref{tab:ch02naobreaks} does that. The point is not
that each specification reaches significance --- two do not --- but that they agree on the year,
including the superseded extract used for the conference presentation that preceded this work.

% GENERATED by diagnostics/build_rc_tables.R - do not hand-edit.
\begin{table}[htbp]\centering
\begin{threeparttable}
\caption{Where the break falls, by specification}
\begin{tabular}{lcccl}
\toprule
Specification & Break & sup-$F$ & Perm.\ $p$ & Bai--Perron \\
\midrule
\textbf{Count, verdict-bearing extracts} (chosen) & 2015 & 16.13$^{***}$ & 0.0002 & 1 break at 2015 \\
Count, all extracts (42-year basis) & 2015 & 12.33$^{**}$ & 0.0010 & 1 break at 2015 \\
Original v2 extract (Manchester spine) & 2015 & 19.05$^{***}$ & 0.0002 & 1 break at 2015 \\
Proportion, verdict-bearing extracts & 2015 & 6.14 & 0.0726 & 1 break at 2015 \\
Proportion, all extracts & 2017 & 7.65$^{*}$ & 0.0274 & 1 break at 2017 \\
Failure within the body's own control & 2015 & 4.02 & 0.2970 & none selected \\
\bottomrule
\end{tabular}
\begin{tablenotes}[flushleft]\footnotesize
\item The question this table answers is where the break falls, not whether it is significant.
\item 5 of the 6 specifications place the break in 2015, including the v2 extract used for the
conference presentation that preceded this work --- a different instrument, applied to a corpus
since expanded by about a tenth.
\item 2 of those 5 cannot distinguish the break from noise at conventional levels; they are shown
with their $p$ values rather than omitted. The agreement is in the date, not in the strength of
the evidence for it.
\item The chosen specification is the count of serious findings estimated from verdict-bearing
extracts. The alternatives are plotted in Figure~\ref{fig:ch02naoalt}.
\item Leave-one-out on the chosen series returns the same break year in 40 of 41 runs.
\item $^{*}p<.05$; $^{**}p<.01$; $^{***}p<.001$, permutation test.
\item Source: \texttt{diagnostics/fig214\_supporting\_tests.R}.
\end{tablenotes}
\end{threeparttable}
\label{tab:ch02naobreaks}
\end{table}


\subsection*{Provenance of the text: three extraction routes, not one}

An independent audit of this chapter's pipeline on 10 August 2026 established that the corpus is
not on a single extractor, and that earlier text here implied otherwise. The position, per row and
recorded in the data:

\begin{itemize}
\item \textbf{354 of 2188 coded reports (16.2 per cent)} carry a \emph{v2 first-prose
fallback} extract, used where no findings anchor could be located. This is the very text the v3
re-code was built to replace, and its incidence is strongly non-stationary --- far commoner in
2000--07 than later. Dropping every fallback row leaves the break year and the direction of the
finding unchanged, and slightly \emph{increases} the measured rise, so the concern that the middle
period is artificially depressed does not hold. The inconsistency is real and is reported here
rather than left implicit.
\item \textbf{47 reports} were read by optical character recognition (Apple Vision), almost all
1984--88, because they were digitised as page images with no text layer. They went through the
same anchor rules and the same prompts. In 1984--95 these rows are coded a serious failure
8.5 per cent of the time against 15.2 per cent for native-text rows --- a difference of
-6.6 points against a two-standard-error band of 9.0, so not statistically distinguishable, on
near-identical sufficiency rates (70.2 against 70.2 per cent). An open question, not a finding.
\item \textbf{Eighteen reports} (1999) were re-extracted with a strengthened guard against
contents-page matches. The guard changed one of the eighteen.
\end{itemize}

A single re-extraction of the whole corpus under one set of rules would remove this untidiness.
It has not been done, and the reader is entitled to know that before weighing the series.

\section{RC09 --- Does the rise reflect government performance, or how the Office writes?}

This is the most serious objection to Figure 2.13, and unlike the others it is internal to the
instrument rather than to the data. It was raised by the author and is tested here directly.

\subsection*{The objection}

The severity scale's threshold for a serious failure is quantified: a cost overrun of roughly
half or more, a majority of a portfolio failing, a qualified audit opinion. A coder can only
apply those triggers to text that states them. RC08 established that this matters --- of 63
reports whose full text codes as a serious failure, a shorter summary version of the same report
agreed on only 28, recording 30 as the milder category, because the summary omitted the
quantified detail rather than the criticism. If National Audit Office reports have become more
quantified across forty years, the coder would record more serious failures without anything
changing in government performance, and the trend in Figure 2.13 would be an artefact of house
style.

\subsection*{The test}

A stratified sample of 416 coded reports (about 110 per period) was re-extracted and the text
measured directly: digit density, explicit percentage mentions, over-run and majority language,
and the verdict vocabulary the scale keys on. Currency mentions were computed and then
\emph{discarded}: the pound sign does not survive text extraction from the older and
image-scanned reports, so the measure was zero for every report before 1997 --- an artefact of
extraction, not an absence of money in the text, and retaining it would have biased the control.

Quantification does rise, though not smoothly. Mean digit density runs 0.0181 in 1984--96,
0.0154 in 1997--2004, 0.0132 in 2005--13 and 0.0264 in 2014--25; explicit percentage mentions run
1.15, 1.60, 1.19 and 3.18. The increase is concentrated in the most recent period rather than
accumulating steadily.

\subsection*{The result}

\textbf{The mechanism is real.} Digit density is a strong and highly significant predictor of a
report being coded a serious failure (+46.2 in log-odds, $p=0.0002$). Reports containing more
numbers are more likely to be recorded as serious, holding the period constant.

\textbf{The trend survives it.} In a logistic model of serious failure on year alone, the
coefficient is +0.439 per decade ($p=0.0003$). Adding digit density, quantification language, verdict
vocabulary and an indicator for whether the extract carried an explicit verdict reduces it to
+0.284 ($p=0.0262$) --- an attenuation of about 35 per cent, with the trend still significant.

\textbf{But the adjusted shape is not the unadjusted shape.} Replacing the linear year term with
period indicators, every period after 1984--96 sits above it --- 1997--2004 +1.13 ($p=0.0096$),
2005--13 +0.81 ($p=0.0581$), 2014--25 +1.03 ($p=0.0124$) --- but the three later periods are not
distinguishable from one another, and the ordering is not monotonic. Once quantification is
controlled, the pattern reads more like a step change after the mid-1990s that then plateaus
than like a steadily compounding rise.

\subsection*{What may be claimed}

That roughly a third of the measured rise is attributable to reports becoming more quantified,
and that the remainder is not. The defensible formulation is a substantial increase in serious
delivery failures identified, of which about a third reflects changes in how the Office writes
rather than in what it found.

Three limits belong with that. The sample is 416 reports, not the full corpus. Digit density is a
crude proxy for quantification and is collinear with the other text measures, so the individual
coefficients should not be read separately. And the period-indicator result sits awkwardly beside
a break dated 2015: the break test is computed on the annual count series, this model on
individual reports with different controls, so they are not in contradiction --- but a reader
will notice, and the tension is better acknowledged than left implicit.

\subsection*{Reproduction}

\texttt{nao\_extension/pdf\_era/text\_features\_sample.py} draws the sample and measures the
text; the models are estimated in
\texttt{nao\_extension/diagnostics/fig214\_supporting\_tests.R}, which writes every figure quoted
above to \texttt{diagnostics/fig214\_computed\_values.csv}.

% GENERATED by diagnostics/build_rc_tables.R - do not hand-edit.
\begin{table}[htbp]\centering
\begin{threeparttable}
\caption{Is the rise a product of reports becoming more quantified?}
\begin{tabular}{lcc}
\toprule
 & Year term & Period terms \\
\midrule
Year (per decade), unadjusted & +0.439$^{***}$ & \\
Year (per decade), adjusted & +0.284$^{*}$ & \\
\quad \emph{attenuation} & \emph{35\%} & \\
\addlinespace
1997--2004 (vs 1984--96) & & +1.13$^{**}$ \\
2005--13 (vs 1984--96) & & +0.81 \\
2014--25 (vs 1984--96) & & +1.03$^{*}$ \\
\addlinespace
Digit density & \multicolumn{2}{c}{+46.2$^{***}$} \\
\bottomrule
\end{tabular}
\begin{tablenotes}[flushleft]\footnotesize
\item Logistic models of whether a report is coded a serious delivery failure. $n=416$ reports,
stratified by period. Adjusted models control for digit density, percentage and over-run language,
verdict vocabulary, and whether the extract carried an explicit verdict. Coefficients are log-odds.
\item \textbf{The mechanism is real and the trend survives it.} Digit density strongly predicts a
serious code, so reports containing more numbers are likelier to be recorded as serious. Adjusting
for it removes about 35 per cent of the year effect and leaves the remainder.
\item \textbf{The adjusted shape differs from the unadjusted shape.} With period terms, all three
post-1996 periods sit above 1984--96 but are not distinguishable from one another: adjusted, the
pattern reads as a step change after the mid-1990s rather than a compounding rise.
\item Currency mentions were measured and discarded: the pound sign does not survive extraction
from older and image-scanned reports, so the measure was zero for every pre-1997 report.
\item $^{*}p<.05$; $^{**}p<.01$; $^{***}p<.001$.
\item Source: \texttt{pdf\_era/text\_features\_sample.py}, models in
\texttt{diagnostics/fig214\_supporting\_tests.R}.
\end{tablenotes}
\end{threeparttable}
\label{tab:ch02naoquant}
\end{table}


\section{RC10 --- government U-turns: three alternative specifications}

Figure 2.15 counts U-turns per year from \emph{The Annual Register}, 1979--2025. The headline
measure is a raw annual count of commitments reversed by the government that made them, dated to
the government decision that reversed them. This check reports three alternative specifications.
One of them reverses the direction of the finding, and the reason it must not be adopted is the
substance of this section.

% GENERATED by data_workspace_2026-07/policy_instability/uturns/gen_rc09_table.py
% Do not edit by hand — re-run the script.

\begin{threeparttable}
\caption{Government U-turns per year under four specifications}
\begin{tabular}{lrrrrr r}
\toprule
Specification & Thatcher & Major & Blair/Brown & Cameron & May to Starmer & $N$ \\
 & 1979--90 & 1991--96 & 1997--09 & 2010--16 & 2017--25 & \\
\midrule
All events (headline) & 5.00 & 7.17 & 3.15 & 3.00 & 6.11 & 220 \\
Per 100,000 characters & 3.87 & 5.87 & 3.56 & 5.46 & 10.15 & 220 \\
High confidence only & 0.75 & 1.50 & 0.92 & 1.29 & 3.33 & 69 \\
Narrow testability only & 3.92 & 5.67 & 2.69 & 2.71 & 5.33 & 183 \\
Manifesto pledges, strict & 0.17 & 0.33 & 0.15 & 0.43 & 0.56 & 14 \\
Manifesto pledges, broad & 0.17 & 0.50 & 0.38 & 0.57 & 1.00 & 23 \\
\bottomrule
\end{tabular}
\begin{tablenotes}\footnotesize
\item Events per year, counted by the year of the government decision that reversed the commitment. $N$ is the total number of events entering each specification.
\item \emph{Per 100,000 characters} divides by the length of the \emph{Annual Register} chapters read for that year. The Register printed 127,637 characters a year on average in 1979--89 against 55,993 in 2010--19, and the reduction is a step rather than a trend.
\item \emph{High confidence only} retains events the coder recorded as high confidence. \textbf{This specification is reported for completeness and should not be used}: 83 per cent of pre-1997 events are coded low or medium confidence against 56 per cent after, so the filter removes early events because they are early. See the discussion.
\item \emph{Narrow testability} retains events whose commitment was narrowly specified, a flag that differs by only 8 percentage points across the two halves of the window.
\item \emph{Manifesto pledges, strict} requires that the manifesto pledge was itself the commitment reversed; \emph{broad} accepts any event whose substance also appeared in a manifesto. \textbf{Eight of the nine events separating them fall after 1997}, so the broad reading roughly doubles the apparent rise. Figure 2.15 plots the strict count.
\end{tablenotes}
\end{threeparttable}

\begin{threeparttable}
\caption{Coding flags by period}
\begin{tabular}{l rr r}
\toprule
Flag & 1979--96 & 1997--2025 & Difference \\
\midrule
Confidence low or medium & 83\% & 56\% & +26pt \\
Judged on one pass only & 17\% & 44\% & -26pt \\
Commitment not stated in source & 40\% & 50\% & -10pt \\
Testability recorded as broad & 21\% & 13\% & +9pt \\
No determinable reversal date & 11\% & 3\% & +8pt \\
Split judging vote & 14\% & 8\% & +6pt \\
Found by one stage-one pass & 17\% & 16\% & +1pt \\
Stage-one duplication flag & 1\% & 0\% & +1pt \\
Short quote unmatchable & 3\% & 3\% & -1pt \\
\bottomrule
\end{tabular}
\begin{tablenotes}\footnotesize
\item Percentages of the 103 events dated before 1997 and the 117 dated from 1997 onwards.
\item The mechanical flags --- stage-one duplication, unmatchable short quotes --- are flat across the window. The flags that skew are those recording the coder's own uncertainty, which is the basis for excluding the high-confidence specification above.
\end{tablenotes}
\end{threeparttable}


\subsection*{Per 100,000 characters: the finding MOVES}

The \emph{Annual Register} printed 127,637 characters a year on average in 1979--89 against 55,993
in 2010--19, a reduction of 2.2 times. Dividing the annual count by chapter length reverses the
early--late comparison: the Thatcher years fall from 5.08 to 3.94 events and the Cameron years rise
from 3.00 to 5.46, so on the normalised measure the later period is the higher one.

The raw count was chosen as the headline before any series had been computed, on the ground that
chapter length is a publisher's page budget rather than a measure of political event-space, and
that the denominator carries its own structural break --- flat to 1999, then two steps down.
Normalising by it would impose the inverse of an editorial decision on the series, and a break test
computed on the normalised measure would partly be detecting when the Register's page allocation
was cut. That argument is unchanged by the direction in which the normalised series happens to run.

Both series are reported. They agree that the period from 2017 is elevated; they disagree about
whether the 1980s and 1990s were also high. That is the boundary between what this measure supports
and what it does not.

\subsection*{High confidence only: the finding REVERSES, and the specification is inadmissible}

Restricting to events the coder recorded as high confidence produces a rising series --- 0.75 events
a year under Thatcher against 3.33 from 2017 --- which matches the expectation registered before the
series existed.

It should not be used. The second table shows why: 83 per cent of events dated before 1997 are coded
low or medium confidence, against 56 per cent of those from 1997 onwards, and 44 per cent of late
events were settled on a single confident judging pass against 17 per cent of early ones. Confidence
is therefore era-dependent by construction. The coder is less certain about 1985 than about 2021
because less is recorded and less is knowable about the earlier cases, not because the earlier events
are worse established as U-turns. Filtering on confidence removes early events \emph{because they are
early}. Using an era-dependent filter to adjudicate an era comparison is circular, and it would
deliver the expected answer by the one route incapable of supporting it.

The specification is reported here rather than omitted. A reader who reconstructs it from the
replication data will find it, and an unexplained omission is worse than a documented rejection.

\subsection*{Narrow testability only: the finding SURVIVES}

The \texttt{testability} flag records whether the reversed commitment was narrowly or broadly
specified. It differs by 8 percentage points between the two halves of the window, against 26 points
for confidence, so it is not contaminated in the same way. Restricting to narrowly specified
commitments retains 184 of 221 events and preserves the direction of the headline series: 4.00 events
a year under Thatcher and 5.67 under Major, against 2.69 and 2.71 in the Blair--Brown and Cameron
periods, rising to 5.33 from 2017.

\subsection*{What the flag table shows in its own right}

The mechanical flags --- stage-one duplication, unmatchable short quotes, re-quoting --- run at
0--3 per cent and are flat across the window. The pipeline is not era-biased. The flags that do skew
all record the coder's own uncertainty: confidence, judging passes, breadth of the commitment, and
whether a reversal date could be determined at all. The early record is harder to code, and the
instrument records that it is.

\subsection*{Reproduction}

\texttt{policy\_instability/uturns/}\allowbreak\texttt{gen\_rc09\_table.py} generates both tables above from
\texttt{coding/uturns\_deduped.csv}; \texttt{build\_fig\_2\_16.R} draws the figure from
\texttt{series/uturns\_series.csv}. The deduplication that produces the 221 events from 250 coded
rows is documented in \texttt{DEDUPLICATION\_RULE.md} and logged case by case in
\texttt{coding/DEDUPE\_LOG.csv}; every post-coding intervention is recorded in
\texttt{POST\_INSTRUMENT\_LOG.md}.

\section{RC11 --- the shape of the turbulence index, and the specifications not chosen}

Figure 2.16 fits a \emph{segmented} model: the index is flat, then bends upward. Every alternative
shape considered is recorded here, with the criteria that chose between them. The question was put
directly --- does the break apparatus used elsewhere in this chapter assume a stepped form that the
data do not support? --- and it was settled on a criterion rather than by argument.

\subsection*{Six shapes, four criteria}

All six models were fitted to the same series (canonical PC1, $n=47$), in the same estimator family.
Parameter counts include the residual variance \emph{and} any searched break date or knot, and BIC
is computed identically for all six as $n\log(\mathrm{RSS}/n)+k\log n$, so the numbers are
comparable. The leave-one-out criterion re-searches the break date or knot \emph{inside each fold},
charging the flexible models honestly for their search; it is the one criterion that does not rest
on a parameter-counting convention.

\begin{center}
\begin{tabular}{lccccc}
\toprule
Model & $k$ & BIC & $\Delta$BIC & LOO-CV RMSE & AR(1) BIC \\
\midrule
Spline (thin-plate, edf 8.33) & 9.33 & 8.00 & \textbf{0.00} & 0.921 & --- \\
Break (level$+$trend, 2015) & 6 & 22.74 & 14.74 & 1.097 & 150.17 \\
\textbf{Segmented (bend, knot 2012)} & 5 & 22.93 & 14.92 & 1.502 & \textbf{148.86} \\
Quadratic & 4 & 23.37 & 15.37 & 1.161 & 149.31 \\
Linear trend & 3 & 41.11 & \textbf{33.11} & 1.443 & 153.95 \\
Flat (no trend) & 2 & 43.35 & \textbf{35.35} & 1.493 & 150.69 \\
\bottomrule
\end{tabular}
\end{center}

Two results, of different strength. \textbf{Flat and linear are decisively rejected} --- $\Delta$BIC
27.4 and 16.8, where a difference above 10 is conventionally decisive, and the ordering is the same
on all four criteria. \textbf{Break, spline and segmented are indistinguishable} --- $\Delta$BIC
0.00, 0.40 and 0.71, all inside the band within which no preference can be claimed. The data
require a late shift and are silent on whether it is a jump, a bend or a smooth curve.

\subsection*{The shape, stated plainly}

The segmented fit puts the slope at $-0.0002$ per year before 2013 and $+0.355$ per year after; the
break fit puts it at $-0.004$ before 2015 with a level jump of $+2.16$. \textbf{There is no gradual
rise.} The pre-shift slope is not merely small, it is indistinguishable from zero in both
interpretable fits. A straight line fitted to the whole series returns a highly significant trend
($t=4.03$, $p<0.0001$), and that statistic is arithmetically correct while describing a shape the
series does not have; it should not be quoted without the shape beside it.

\subsection*{Three controls}

\textbf{Not an artefact of the PCA weights.} PC1's loadings are estimated on the full sample, which
can manufacture late structure. The unweighted mean of the six $z$-scores, with no estimated weights
at all, correlates $r=0.989$ with PC1 and returns the same verdict: linear 11.24 worse and flat
25.85 worse than the leading models, with a segmented knot at 2012.

\textbf{All six components participate.} Mean $z$ before and after 2013: rebellions $+1.245$,
dislocation $+1.182$, electoral volatility $+1.095$, PM churn $+0.577$, departmental churn $+0.500$,
U-turns $+0.277$. Six of six move upward. The components are correlated, so this is description
rather than an independent sign test, but it establishes that the shift is not one or two series
dragging an index.

\textbf{The step does not survive the end date.} Ending the series in 2019, before the pandemic,
reverses the ordering: spline 0.00, segmented 2.98, break 4.87, quadratic 8.79, flat 10.50, linear
10.71, with the knot at 2013. Flat and linear remain rejected; the smooth alternatives now beat the
step. The existence of the late shift is robust to the end date. The claim that it is a
discontinuity is not.

\subsection*{The components do not individually support breaks}

Running the same selection exercise on each component separately, only one of six prefers a break
model. Departmental churn prefers a plain linear trend ($\Delta$BIC 6.51 against the break),
consistent with its Poisson sup-LR result of $p=0.10$; PM churn and U-turns prefer \emph{flat};
electoral volatility and rebellions prefer a spline, and rebellions rejects a break by $\Delta$BIC
31.75. Only parliamentary dislocation prefers one. This is consistent with the previous control
rather than against it: individual annual series are too noisy to locate a shift, while all six
lean the same way, so aggregation recovers a signal none of them shows alone. It does mean that
\textbf{the shift belongs to the index and not to its parts}, and that break dates should not be
claimed for individual component series on this evidence.

\subsection*{Where the turn falls, and how firmly}

A block bootstrap (1{,}000 draws, block length 4) puts the knot's 90\% interval at
\textbf{2007--2017}, with a median of 2013 and only 20.1 per cent of draws on the modal year;
$P(\text{knot} \in 2011\text{--}2015) = 0.705$. Leaving out each component in turn moves the knot to
2012, 2012, 2013, 2013 and 2014 --- except electoral volatility, whose removal moves it to
\textbf{2004} and roughly halves the post-shift slope. \textbf{The existence of the shift rests on
all six components; its date rests disproportionately on one.} The main text therefore dates the
turn as a range in the early-to-mid 2010s and Figure 2.16 draws it as a band rather than a line.

\subsection*{Relation to RC06}

RC06 shows that a single-break test applied to a curved series returns a break wherever the best
split happens to land, with a bootstrap interval spanning almost the whole window --- and warns that
a wide interval is a symptom of misspecification. The turbulence index has a wide interval, so the
warning applies and was tested here. The outcome differs: for this series the wide interval reflects
uncertainty about \emph{when}, not evidence against a shift, because the flat and linear models are
rejected outright and three different late-shift shapes all fit. The lesson RC06 draws --- treat the
shape as the finding and any single break year as conditional --- is adopted here in full.

\subsection*{Reproduction}

\texttt{nonlinear\_modelling/horse\_race.R} runs the six-model comparison and the per-component
races; \texttt{control\_equal\_weight.R} runs the equal-weight and end-date controls;
\texttt{knot\_robustness.R} runs the leave-one-out and bootstrap on the knot;
\texttt{shape.R} reports the fitted slopes. All are local R, seed 42, and write to
\texttt{nonlinear\_modelling/}. The reasoning and the result are recorded in
\texttt{NONLINEAR\_MODEL\_SELECTION.md}, including the fact that the exercise \textbf{overturned the
expectation it was set up to test}: a smooth trend was argued for in advance and the criteria
rejected it.

\section{RC12 --- the volatility curse, and four claims tested and not found}

Chapter 1 names Campello and Zucco's volatility curse as the mechanism that would carry the book's
argument from measurement to consequence: where outcomes are driven by forces beyond a government's
control, voters reward and punish incumbents for conditions they did not create, and elections stop
transmitting a democratic signal. If that mechanism operates in Britain, the chapter's own data
should show it. This section reports the test, which is direct, and the result, which is that the
mechanism is not found. Four further claims a reader might expect the chapter to make were tested
in the same exercise and none of them survived either.

\subsection*{Where the instrument comes from}

Figure 2.9 already fits, for each item on the public agenda, a model of public worry on the
real-world condition that item names --- worry about inflation on the inflation rate, worry about
unemployment on the unemployment rate, and so on. \textbf{The fit of that model is the strength of
the attribution link}: how closely public concern tracks the world. Nothing further needs to be
built. If the volatility curse is operating, that tracking should weaken over the window. Eleven
items have at least thirty usable years. The test uses no data from Chapters 3 to 9, so it cannot
contaminate the suspects.

\subsection*{The result depends entirely on which measure is used}

% GENERATED by data_workspace_2026-07/_exploratory_ch3-10/attribution_test/rc12_nulls.R
% Do not edit by hand — re-run the script.

\begin{threeparttable}
\caption{The volatility curse: two measures of the attribution link, by era}
\begin{tabular}{lrrr r rrr}
\toprule
 & \multicolumn{3}{c}{$R^2$} & & \multicolumn{3}{c}{Standardised slope $|\beta|$} \\
\cmidrule(lr){2-4}\cmidrule(lr){6-8}
Item & pre-2000 & 2000--14 & post-2015 & & pre-2000 & 2000--14 & post-2015 \\
\midrule
Crime & 0.623 & 0.051 & 0.275 & & 1.053 & 0.495 & 1.535 \\
Drugs & 0.877 & 0.110 & 0.439 & & 1.878 & 0.612 & 1.303 \\
Economy (general) & 0.092 & 0.322 & 0.002 & & 0.307 & 1.000 & 0.011 \\
Housing & 0.052 & 0.306 & 0.234 & & 0.324 & 0.544 & 2.375 \\
Immigration & 0.002 & 0.768 & 0.099 & & 0.026 & 2.240 & 0.214 \\
Inflation & 0.904 & 0.064 & 0.797 & & 1.005 & 0.255 & 0.818 \\
Pensions & 0.726 & 0.062 & 0.553 & & 1.811 & 0.233 & 1.668 \\
Petrol / fuel & --- & 0.672 & 0.579 & & --- & 1.056 & 1.214 \\
Poverty & --- & 0.512 & 0.471 & & --- & 0.514 & 0.428 \\
Trade unions / strikes & 0.683 & 0.223 & 0.457 & & 0.845 & 0.291 & 0.258 \\
Unemployment & 0.861 & 0.822 & 0.521 & & 0.896 & 0.692 & 0.632 \\
\midrule
\textbf{Mean} & \textbf{0.536} & \textbf{0.356} & \textbf{0.402} & & \textbf{0.905} & \textbf{0.721} & \textbf{0.950} \\
Median & 0.683 & 0.306 & 0.457 & & 0.896 & 0.544 & 0.818 \\
\bottomrule
\end{tabular}
\begin{tablenotes}\footnotesize
\item Each cell fits public worry on the real-world condition the item names, within the era. Items with fewer than 30 usable years, or fewer than six years in an era, are not reported.
\item The standardised slope multiplies the raw slope by $\mathrm{sd}(\text{condition})/\mathrm{sd}(\text{worry})$ computed on the \emph{full} series, not within the era. Standardising within era returns the correlation coefficient, whose square is the $R^2$ beside it; a common yardstick is what makes the two columns independent tests rather than one test written twice.
\item Paired Wilcoxon, post-2015 against pre-2000: $R^2$ $V=7.0$, $p=0.074$; slope $V=18.0$, $p=0.652$.
\end{tablenotes}
\end{threeparttable}

\begin{threeparttable}
\caption{Attribution and delivery failure: every naive result dies under a time trend}
\begin{tabular}{lrr rr rr}
\toprule
 & \multicolumn{2}{c}{Naive} & \multicolumn{2}{c}{$+$ linear trend} & \multicolumn{2}{c}{Differenced} \\
\cmidrule(lr){2-3}\cmidrule(lr){4-5}\cmidrule(lr){6-7}
Lag & $b$ & $p$ & $b$ & $p$ & $b$ & $p$ \\
\midrule
0 & -46.97 & 0.0001 & -12.42 & 0.3463 & -20.43 & 0.0702 \\
1 & -39.00 & 0.0001 & +2.06 & 0.8787 & +18.74 & 0.2969 \\
2 & -40.22 & 0.0000 & -7.26 & 0.6051 & +11.52 & 0.5359 \\
\bottomrule
\end{tabular}
\begin{tablenotes}\footnotesize
\item Dependent variable is the National Audit Office delivery-failure count; the regressor is the pooled rolling $R^2$ of the attribution models, lagged. Newey--West standard errors.
\item The volatility curse predicts a \emph{negative} coefficient: where the public can no longer attribute outcomes to government, the discipline on government weakens and delivery worsens.
\item Raw correlation between attribution and delivery failure is $r=-0.659$; after removing a linear trend from both it is $r=-0.158$. The naive result is co-trending.
\end{tablenotes}
\end{threeparttable}

\begin{threeparttable}
\caption{Four claims tested and not found}
\begin{tabular}{p{4.4cm}p{4.4cm}l}
\toprule
Claim & Test & Result \\
\midrule
The accountability link has decayed (Campello \& Zucco) & Standardised slope, 9 items, three eras & Not found ($p=0.65$) \\
Attribution decay accompanies turbulence & Pooled rolling $R^2$ against the index & Not found ($\rho=-0.19$, $p=0.29$) \\
Weak attribution predicts delivery failure & Lags 0--2, trend and differenced & Not found (detrended $r=-0.16$) \\
Critical slowing down before the turn & Rolling variance and AR(1) & Not found (variance \emph{falls}, $\rho=-0.10$) \\
\bottomrule
\end{tabular}
\begin{tablenotes}\footnotesize
\item A fifth claim --- that turbulence precedes policy failure --- is tested in \texttt{turbulence\_to\_failure.R} and is likewise not found: naive $b=+2.58$, $p=0.003$, falling to $b=-0.19$, $p=0.73$ with a trend, and Granger-null in both directions.
\item Against these, one positive finding from the same exercise. Mean pairwise correlation between the six dimensions falls from $+0.445$ (1979--89) to $+0.042$ (2015--25), $\rho=-0.434$, $p=0.007$: the dimensions are becoming \emph{more} independent, which is evidence for a formative reading of the index and the reason PC1 explains only part of the variance.
\end{tablenotes}
\end{threeparttable}



On $R^2$ the curse appears to be confirmed. Mean $R^2$ falls from 0.536 before 2000 to 0.402 after
2015, it falls in seven of the nine testable items, and the pooled rolling $R^2$ declines steeply
across the window ($\rho=-0.730$, $p<0.0001$). Reported alone, that is a publishable finding.

\textbf{It does not survive the obvious check.} $R^2$ measures the \emph{proportion} of variance
explained, so it falls mechanically whenever the condition being tracked stops varying --- even if
the public tracks that condition exactly as well as before. This is range restriction, and it is not
a decay of attribution. The measure that is immune to it is the slope, scaled by a yardstick held
constant across eras. On the standardised slope the mean does not fall at all (0.905 to 0.950), the
median falls slightly (0.896 to 0.818), the count of declining items drops from seven to six, and
the paired test returns $p=0.65$. \textbf{On the measure that range restriction cannot fake, the
decay is absent.}

\subsection*{The middle era is the anomaly, not the present}

The three-era columns show the shape the two-era comparison hides. Mean $R^2$ runs 0.536, 0.356,
0.402: a dip and a partial recovery, not a decline. Inflation is the clearest case. Its $R^2$
collapses from 0.904 before 2000 to 0.064 in 2000--14 and returns to 0.797 after 2015, and the
standard deviation of inflation itself over those eras is 4.14, 1.36 and 3.35. \textbf{The collapse
is the Great Moderation.} There was almost no inflation to attribute, so the fit vanished; when
inflation returned after 2015, attribution returned with it, close to its 1980s strength. Trade
unions follow the same pattern for the same reason.

\textbf{The explanation does not extend to every item, and should not be claimed to.} Pensions shows
the same V ($R^2$ 0.726, 0.062, 0.553) but its condition was \emph{most} variable in the middle era,
not least (standard deviations 1.08, 3.38, 1.50), so range restriction cannot be what produced its
dip. The common finding across items is the V-shape and the absence of a monotone decline; the
range-restriction account explains why the V occurs in two of the three cases examined, and the
third is unexplained.

\subsection*{Attribution does not move with the index}

If turbulence severs the accountability link, the two should move together. The pooled rolling $R^2$
and the turbulence index correlate at $r=-0.274$, Spearman $\rho=-0.192$, $p=0.29$ over 32 years.
The sign is the predicted one and the magnitude is not distinguishable from zero. Even had the decay
been real, it does not track turbulence year to year.

\subsection*{The output half fails on the same pattern}

The curse has a second limb: if the accountability link is severed, the discipline on government
weakens and delivery gets worse. The National Audit Office series is the chapter's only output
measure, and the second table above regresses it on lagged attribution. Every naive coefficient is
large, correctly signed and highly significant. \textbf{Every one of them dies when a linear trend is
added}, and at lag 1 the sign flips. The raw correlation is $r=-0.659$; detrended it is $r=-0.158$.
One specification of nine survives at the ten per cent level. Nothing may be built on it.

\subsection*{No early warning}

A system approaching a critical transition is expected to show rising variance and rising
autocorrelation beforehand --- critical slowing down. On eleven-year rolling windows of the
detrended index, variance does not rise ($\rho=-0.10$, $p=0.56$) and autocorrelation does not rise
($\rho=+0.12$, $p=0.46$). The result is unchanged on windows of eight, ten and fifteen years and on
the raw as well as the detrended series. \textbf{The turn of the mid-2010s was not announced in
advance by the statistical signature of an approaching tipping point}, and tipping-point language
should not be used of it.

\subsection*{The pattern across all five tests is itself the finding}

Five causal claims have now been tested on these data: attribution decay, attribution against the
index, attribution to delivery failure, turbulence to delivery failure, and critical slowing down.
\textbf{In four of the five the naive result is strong, correctly signed and significant, and in
every one of those four it is entirely accounted for by common trend.} Four decades of British
political series trend together, and a naive analysis of almost any pair of them will produce a
mechanism that looks publishable and is not there. Given how much commentary about British political
decline rests on exactly such pairings, this is worth reporting as a result rather than buried as a
failure.

\subsection*{What the chapter may still claim, and what it may not}

\textbf{The chapter may not claim that the accountability link has decayed.} Its own data will not
carry it, and the check that overturns it is one a competent referee will run in an afternoon.
\textbf{It may report the volatility curse as a mechanism tested directly and not found} --- which
is a contribution, since the mechanism is routinely asserted of Britain and, so far as we are aware,
has not previously been tested on four decades of British agenda data.

\textbf{One consequence must be stated plainly.} The literature audit identified the volatility
curse as the argument that rescues the composite from the objection that it measures turbulence in
the machinery of government rather than democratic decline. That bridge does not survive testing, so
\textbf{the construct-validity question is left open, and the honest framing is the narrow one}: the
index measures turbulence in the machinery of government, its rise after the mid-2010s is large and
robust, and whether that constitutes democratic decline is a question the measurement cannot settle.

\textbf{Against the five nulls, one positive finding.} Mean pairwise correlation between the six
component dimensions falls from $+0.445$ in 1979--89 to $+0.042$ in 2015--25 ($\rho=-0.434$,
$p=0.007$). The dimensions are becoming \emph{more} independent over time, not less, so whatever is
happening is not contagion between them. This supports a \textbf{formative} reading of the index ---
the six things constitute turbulence rather than reflecting a common latent cause --- and it is also
the answer to the natural objection that PC1 explains only 36.4 per cent of the variance. On a
formative index that figure is not a defect to be minimised; a reflective index of six measures of
one underlying thing would be expected to do considerably better, and the fact that this one does
not is evidence about what kind of object it is.

\subsection*{Reproduction}

\texttt{attribution\_test/rc12\_nulls.R} runs all five tests reported here and writes this section's
tables; \texttt{attribution\_test/run\_attribution.R} runs the original $R^2$ and rolling-window
analysis; \texttt{nonlinear\_modelling/turbulence\_to\_failure.R} runs the turbulence-to-failure
test. All are local R, seed 42, no external calls. Outputs are written to
\texttt{attribution\_test/} as \texttt{RC12\_NULLS\_RESULTS.txt} and five CSVs. The reasoning,
including the numbers that an earlier unscripted run reported and that this script does not
reproduce, is recorded in \texttt{attribution\_test/ATTRIBUTION\_NOTES.md}.
